\documentclass[11pt]{article}
\usepackage[margin=1in]{geometry}
\usepackage{enumitem}
% =============================================================================
% Preambles/header.tex — shared LaTeX/Beamer preamble for this project
%
% Use from a Beamer lecture by prepending:
%     \input{header}
% and compiling with TEXINPUTS=../Preambles:$TEXINPUTS (the /compile-latex
% skill does this automatically).
%
% PALETTE CONTRACT. Color names defined here MUST match the names in
% Quarto/theme-template.scss so Beamer and Quarto renderings use the same
% palette. The scripts/check-palette-sync.sh script enforces this.
%
% To customize: edit the HEX values below AND the matching $variable in
% Quarto/theme-template.scss, then run scripts/check-palette-sync.sh.
% =============================================================================

% -----------------------------------------------------------------------------
% Engine detection + encoding
%
% Our /compile-latex skill uses XeLaTeX, where inputenc is unnecessary and
% can produce encoding warnings. Only load inputenc under pdfTeX.
% -----------------------------------------------------------------------------
\usepackage{iftex}
\ifPDFTeX
  \usepackage[utf8]{inputenc}
\fi

\usepackage{amsmath, amssymb, amsthm}
\usepackage{booktabs}
\usepackage{graphicx}

% -----------------------------------------------------------------------------
% PALETTE — mirrors Quarto/theme-template.scss variable names.
% Load xcolor and define colors BEFORE anything that references them
% (notably hyperref, hypersetup, and the Beamer color assignments).
% -----------------------------------------------------------------------------
\usepackage{xcolor}

\definecolor{primary-blue}{HTML}{012169}      % headings, accents
\definecolor{primary-gold}{HTML}{B9975B}      % emphasis, borders
\definecolor{highlight-yellow}{HTML}{F2A900}  % markers, alerts
\definecolor{light-bg}{HTML}{E8EDF5}          % title / section backgrounds
\definecolor{jet}{HTML}{1A1A1A}               % body text

% Semantic colors (used in diagrams and annotations)
\definecolor{positive}{HTML}{15803D}          % good / observed / identified
\definecolor{negative}{HTML}{B91C1C}          % bad / problematic / bias
\definecolor{neutral}{HTML}{525252}           % reference / context

% Highlight accents (match .hi-* classes in the SCSS)
\definecolor{hi-slate}{HTML}{314F4F}
\definecolor{hi-green}{HTML}{2E8B57}
\definecolor{hi-red}{HTML}{C41E3A}

% Semantic aliases so skill templates and snippets can write
% \textcolor{accent}{...} without hard-coding "primary-blue".
\colorlet{accent}{primary-blue}
\colorlet{accent2}{primary-gold}

% -----------------------------------------------------------------------------
% Hyperref — loaded AFTER xcolor + palette so link colors resolve.
% -----------------------------------------------------------------------------
\usepackage{hyperref}
\hypersetup{colorlinks=true, linkcolor=primary-blue, urlcolor=primary-blue,
            citecolor=primary-blue, pdfborder={0 0 0}}

% -----------------------------------------------------------------------------
% Course code — set once per deck (after \input{header}, before \begin{document})
% via \coursecode{CS401}. Shown in the footer on every slide so a deck's course
% is identifiable without opening the title page. Empty by default.
% -----------------------------------------------------------------------------
\makeatletter
\newcommand{\coursecode}[1]{\gdef\@coursecodetext{#1}}
\gdef\@coursecodetext{}
\makeatother

% -----------------------------------------------------------------------------
% Beamer theme assignments (only apply under Beamer).
%
% We detect Beamer via \@ifclassloaded{beamer}, which is the documented,
% non-internal way. This must be wrapped in \makeatletter/\makeatother
% because \@ifclassloaded contains an @-catcode character.
% -----------------------------------------------------------------------------
\makeatletter
\@ifclassloaded{beamer}{%
  \usefonttheme{professionalfonts}
  \setbeamercolor{structure}{fg=primary-blue}
  \setbeamercolor{frametitle}{fg=primary-blue}
  \setbeamercolor{title}{fg=primary-blue}
  \setbeamercolor{subtitle}{fg=primary-gold}
  \setbeamercolor{author}{fg=primary-blue}
  \setbeamercolor{institute}{fg=neutral}
  \setbeamercolor{date}{fg=primary-gold}
  \setbeamercolor{itemize item}{fg=primary-blue}
  \setbeamercolor{itemize subitem}{fg=primary-gold}
  \setbeamercolor{alerted text}{fg=primary-gold}
  \setbeamercolor{block title}{fg=white, bg=primary-blue}
  \setbeamercolor{block body}{bg=light-bg}
  % Semantic block variants: block=definition/concept (blue), exampleblock=
  % worked example (green/positive), alertblock=key takeaway or caution (gold).
  % Keep to <=2 colored boxes per slide across ALL three variants (INV-7).
  \setbeamercolor{block title example}{fg=white, bg=positive}
  \setbeamercolor{block body example}{bg=light-bg}
  \setbeamercolor{block title alerted}{fg=white, bg=primary-gold}
  \setbeamercolor{block body alerted}{bg=light-bg}

  % Minimal footer: course code (left) + page X / N (right), no navigation clutter
  \setbeamertemplate{navigation symbols}{}
  \setbeamertemplate{footline}{%
    \color{neutral}\footnotesize\hspace{0.7em}\@coursecodetext\hfill
    \insertframenumber/\inserttotalframenumber\hspace{0.7em}\vspace{0.3em}}
}{}
\makeatother

% -----------------------------------------------------------------------------
% TikZ setup — libraries the snippet gallery uses
% -----------------------------------------------------------------------------
\usepackage{tikz}
\usetikzlibrary{arrows.meta, positioning, calc, decorations.pathreplacing,
                fit, shapes.geometric, backgrounds}

% Shared TikZ styles — snippets and hand-written diagrams can pick these up
% via `\begin{tikzpicture}[every node/.style={...}]` or by naming specific
% styles (e.g., `\node[dag-node] (X) at (0,0) {X};`).
\tikzset{
  % Node styles for causal DAGs and similar
  dag-node/.style = {
    draw=primary-blue, fill=light-bg, rounded corners=2pt,
    minimum width=1.2cm, minimum height=0.9cm, align=center, font=\small
  },
  decision-node/.style = {
    draw=primary-gold, fill=white, diamond, aspect=1.8,
    minimum width=1.8cm, minimum height=1.2cm, align=center, font=\small,
    inner sep=1pt
  },
  % Edge styles — observed vs counterfactual
  observed-edge/.style = {-{Latex[length=2.5mm]}, thick, draw=jet},
  counterfactual-edge/.style = {-{Latex[length=2.5mm]}, thick, draw=neutral, dashed},
  confound-edge/.style = {-{Latex[length=2.5mm]}, thick, draw=negative, dashed},
  % Data-point styles for plots
  observed-dot/.style = {fill=primary-blue, draw=primary-blue, circle, inner sep=1.2pt},
  counterfactual-dot/.style = {fill=white, draw=neutral, circle, inner sep=1.2pt}
}

% -----------------------------------------------------------------------------
% Convenience macros
% -----------------------------------------------------------------------------
\newcommand{\muted}[1]{\textcolor{neutral}{#1}}
\newcommand{\key}[1]{\textcolor{primary-gold}{\textbf{#1}}}
\newcommand{\good}[1]{\textcolor{positive}{#1}}
\newcommand{\bad}[1]{\textcolor{negative}{#1}}

% Standout frame macro: full-bleed dark background for section transitions.
% Beamer-only (uses \begin{frame}/\setbeamercolor/\usebeamerfont) -- do not
% call from an article-class file (e.g. Notes/<CODE>/*-notes.tex). \muted,
% \key, \good, \bad, and \coursecode above are plain xcolor/gdef and are
% safe to use from either class; verified when Notes/article-class support
% was added.
% Expand: \transitionslide{Your transition title}
\newcommand{\transitionslide}[1]{%
  \begingroup
  \setbeamercolor{background canvas}{bg=primary-blue}
  \begin{frame}[plain]
    \centering
    \vfill
    {\usebeamerfont{title}\color{white}\Large #1\par}
    \vfill
  \end{frame}
  \endgroup
}

% Section divider frame (Beamer-only), an alternative to \transitionslide with a
% darker two-line style: near-black (jet) background, a small white label line
% above a larger gold title, and a thin gold rule along the bottom edge.
% Expand: \sectiondivider{Section 2}{Pointers}
\newcommand{\sectiondivider}[2]{%
  \begingroup
  \setbeamercolor{background canvas}{bg=jet}
  \begin{frame}[plain]
    \vspace*{3.0cm}
    {\color{white}\ttfamily\large #1\par}
    \vspace{0.5cm}
    {\color{primary-gold}\LARGE #2\par}
    \begin{tikzpicture}[remember picture, overlay]
      \draw[primary-gold, line width=2.5pt]
        ($(current page.south west)+(0,0.1)$) -- ($(current page.south east)+(0,0.1)$);
    \end{tikzpicture}
  \end{frame}
  \endgroup
}

% -----------------------------------------------------------------------------
% Channel watermark — "Concepts That Click" (YouTube).
%
% Article-class files (CompetitiveExam Q/A sets, Notes, Assignments) get a
% light diagonal draft watermark on every page plus a centered footer naming
% the channel. Beamer decks get a subtle TikZ background watermark instead
% (the footline already carries the course code). Centralized here so every
% MCQ PDF is branded without per-file edits.
% -----------------------------------------------------------------------------
\makeatletter
\@ifclassloaded{article}{%
  \usepackage{draftwatermark}
  \SetWatermarkText{Concepts That Click}
  \SetWatermarkScale{1}
  \SetWatermarkFontSize{2cm}
  \SetWatermarkLightness{0.86}
  \SetWatermarkAngle{45}
  \usepackage{fancyhdr}
  \pagestyle{fancy}
  \fancyhf{}
  \fancyfoot[C]{\footnotesize\textcolor{neutral}{Concepts That Click~|~YouTube: \href{https://www.youtube.com/@concepts.thatclick}{@concepts.thatclick}}}
  \renewcommand{\headrulewidth}{0pt}
  \renewcommand{\footrulewidth}{0pt}
  \fancypagestyle{plain}{%
    \fancyhf{}%
    \fancyfoot[C]{\footnotesize\textcolor{neutral}{Concepts That Click~|~YouTube: \href{https://www.youtube.com/@concepts.thatclick}{@concepts.thatclick}}}%
    \renewcommand{\headrulewidth}{0pt}%
    \renewcommand{\footrulewidth}{0pt}%
  }
}{}
\@ifclassloaded{beamer}{%
  \setbeamertemplate{background}{%
    \begin{tikzpicture}[remember picture, overlay]
      \node[rotate=45, opacity=0.055, align=center]
        at (current page.center)
        {\Huge\textcolor{neutral}{Concepts That Click}};
    \end{tikzpicture}%
  }
}{}
\makeatother 
\coursecode{JKSSB}
\renewcommand{\thesection}{35.\arabic{section}}
\newtheorem{example}{Example}[section]
\newenvironment{definitionbox}[1]{%
  \par\noindent\textbf{Definition (#1).}\ \itshape
}{\par\vspace{0.3em}}
\title{URL, DNS, HTTP and HTTPS --- Detailed Lecture Notes}
\author{JKSSB Computer Awareness}
\date{\today}

\begin{document}
\maketitle

\section{What a URL is}
A URL (Uniform Resource Locator) is the full address of one resource on the
Internet: one page, file, image, or service. Typing a URL tells the browser
both how to connect and exactly what to fetch.

\begin{definitionbox}{URL}
The complete address of a single Internet resource, naming the protocol
(scheme), the host, and the location of the resource on that host,
including the optional port, path, and query.
\end{definitionbox}

\begin{center}
\begin{tabular}{p{0.24\textwidth}p{0.66\textwidth}}
\toprule
\textbf{Term} & \textbf{Meaning in this lesson} \\
\midrule
URL & The full address of one resource, such as a page or file. \\
Scheme & The protocol word before the colon, such as https or http. \\
Domain name & The readable host name, such as www.example.in. \\
TLD & Top-Level Domain: the last label, such as .in or .com. \\
Subdomain & A prefix section of the name, such as www or mail. \\
Port & The door number on the server (80 for HTTP, 443 for HTTPS). \\
Path & The location of the page inside the site. \\
Query & Parameters after a ?, such as ?topic=dns. \\
IP address & The numeric address of a device (IPv4 is 32-bit, IPv6 is 128-bit). \\
DNS & Domain Name System: translates domain names into IP addresses. \\
HTTP & Hypertext Transfer Protocol: carries web pages. \\
HTTPS & Hypertext Transfer Protocol Secure: HTTP protected by TLS. \\
TLS & Transport Layer Security: encrypts the exchange (successor of SSL). \\
\bottomrule
\end{tabular}
\end{center}

The rest of this lesson uses these terms in one consistent sense. A
\textbf{URL} is the whole address; the \textbf{domain name} is only the host
part inside it. \textbf{DNS} translates that host name into an \textbf{IP
address}. \textbf{HTTP} moves the page, and \textbf{HTTPS} is HTTP carried
inside a \textbf{TLS}-encrypted channel.

\section{Anatomy of a URL}
Consider the example used throughout this lesson:
\begin{center}
\texttt{https://www.example.in:443/study/computer?topic=dns\&lang=en}
\end{center}
Read it left to right. The \textbf{scheme} is \texttt{https}: the protocol
the browser must speak. The \texttt{://} is only a separator between the
scheme and the host. The host is \texttt{www.example.in}, in which
\texttt{www} is the \textbf{subdomain}, \texttt{example} is the registered
domain, and \texttt{.in} is the \textbf{TLD} (Top-Level Domain). Common TLDs
include \texttt{.com}, \texttt{.org}, \texttt{.edu}, \texttt{.gov}, and
country labels such as \texttt{.in}.

After the host comes the optional \textbf{port}, \texttt{:443}: the door
number on the server. The \textbf{path}, \texttt{/study/computer}, locates
the page inside the site. The \textbf{query}, \texttt{?topic=dns\&lang=en},
passes extra parameters to the server. Port and query are optional: when the
port is missing, the browser fills in the default for the scheme.

\begin{definitionbox}{TLD (Top-Level Domain)}
The last label of a domain name, such as .com, .org, .edu, .gov, or a
country label such as .in.
\end{definitionbox}

\begin{example}
Split \texttt{mail.example.com}. The subdomain is \texttt{mail}, the
registered domain is \texttt{example}, and the TLD is \texttt{.com}. One
registered domain can carry many such subdomains (\texttt{www},
\texttt{mail}, \texttt{blog}), each pointing to a different service, while
the domain itself stays the same.
\end{example}

\section{Ports, paths, and queries}
The \textbf{port} selects which service on the server receives the
connection. Port \textbf{80} is the default for HTTP and port \textbf{443}
is the default for HTTPS, which is why both numbers are usually hidden in
the address bar. The \textbf{path} names the page or file within the site,
and the \textbf{query} carries name--value pairs after a \texttt{?} mark.
Neither the port nor the query changes which host is contacted; the host is
fixed by the domain name alone.

\section{IP addresses: the real destination}
An IP (Internet Protocol) address is the numeric address every device uses
to send and receive data on the network. People type names; machines route
by numbers.

\begin{definitionbox}{IP address}
The numeric address of a device on a TCP/IP network, used for routing data
to that device.
\end{definitionbox}

\textbf{IPv4} (Internet Protocol version 4) is \textbf{32-bit} and is
written as four decimal groups, for example \texttt{192.0.2.44}.
\textbf{IPv6} (Internet Protocol version 6) is \textbf{128-bit} and is
written in hexadecimal groups. It is never 256-bit; that figure is a
repeated wrong option. IPv6 exists because the IPv4 space is too small for
the number of connected devices.

\section{DNS: the phone book of the Internet}
DNS (Domain Name System) is the distributed naming system that translates
domain names into IP addresses. It does not carry page content and does not
assign addresses to devices; it only answers the question ``which IP address
belongs to this name?''

\begin{definitionbox}{DNS}
The hierarchical, distributed system that maps each domain name to the IP
address of its server, with answers cached to speed up repeat visits.
\end{definitionbox}

Resolution runs as a short chain. The browser first checks its own cache
for the name. If the name is unknown, the \textbf{resolver} (usually run by
the ISP --- Internet Service Provider) asks the DNS hierarchy: the root
servers, then the TLD servers for labels such as \texttt{.in}, then the
domain's own authoritative server. That authoritative server returns the IP
address, the resolver caches the answer for a fixed time, and the browser
receives the number it can connect to.

\begin{example}
On the office desktop, typing \texttt{www.example.in} first triggers a DNS
lookup. The resolver walks the hierarchy, learns the address (say
\texttt{192.0.2.44}), caches it, and only then does the browser open a
connection to that number. Typing the IP address directly would skip DNS,
but nobody memorises numbers for every site.
\end{example}

\section{HTTP: carrying web pages}
HTTP (Hypertext Transfer Protocol) is the request--response protocol of the
web. The browser sends a request naming the path (and query); the server
returns the page. HTTP runs by default over port 80. Plain HTTP provides no
encryption: the request and the page travel in readable form, so anyone on
the path between browser and server can read them.

\begin{definitionbox}{HTTP}
The protocol by which a web client requests a resource and the server
returns it. Plain HTTP does not encrypt the exchange.
\end{definitionbox}

\section{HTTPS and TLS: the secure form}
HTTPS (Hypertext Transfer Protocol Secure) is HTTP protected by TLS. TLS
(Transport Layer Security) is the successor of the older SSL (Secure
Sockets Layer); the phrase ``SSL/TLS'' survives only as a name for the same
family of protection. Before any page moves, the browser and server run a
TLS handshake: the server proves its identity with a certificate, and the
two sides agree on encryption keys. The HTTP request and response then
travel inside that encrypted channel.

\begin{definitionbox}{HTTPS}
HTTP carried inside a TLS-encrypted channel, by default over port 443 and
shown with a padlock in the address bar.
\end{definitionbox}

TLS provides three things at once: \textbf{confidentiality} through
encryption (only the two ends can read the data), \textbf{integrity}
(tampering is detected), and server \textbf{authentication} (the certificate
ties the site to its claimed name). HTTPS therefore protects the trip; it
is not a different page language.

\section{The full flow: from typed URL to secure page}
Put the pieces together for \texttt{https://www.example.in/study/computer}.
The browser parses the URL: scheme \texttt{https}, host
\texttt{www.example.in}, path \texttt{/study/computer}. DNS resolves the
host to an IP address. The browser opens a connection to that IP on port
443, the HTTPS default. The TLS handshake secures the channel, and the HTTP
request for the path travels inside it. The server returns the page, and the
browser renders it with the padlock shown.

\section{Exam retrieval and common confusions}
Six distinctions prevent most errors on this topic.
\begin{enumerate}
  \item A \textbf{URL} addresses one resource; the \textbf{domain name} is
    only the host part inside it.
  \item \textbf{DNS} translates names into IP addresses. It does not assign
    addresses and does not carry page content.
  \item \textbf{IPv4} is 32-bit; \textbf{IPv6} is 128-bit, never 256-bit.
  \item \textbf{HTTP} uses port 80; \textbf{HTTPS} uses port 443. Do not
    reverse them.
  \item \textbf{HTTPS} is HTTP over \textbf{TLS}. HTTP alone gives no
    encryption.
  \item The \textbf{scheme} is the protocol word (\texttt{https}), not the
    domain name that follows it.
\end{enumerate}

\begin{example}
An item asks which part of \texttt{https://mail.example.com/inbox} is the
subdomain. Trace the host \texttt{mail.example.com} from the right: the TLD
is \texttt{.com}, the domain is \texttt{example}, and the remaining prefix
\texttt{mail} is the subdomain. If the stem instead asks for the scheme,
the answer is \texttt{https}; if it asks for the default port of that
scheme, the answer is 443.
\end{example}

\subsection*{Exercises}
\begin{enumerate}[label=\arabic*.]
  \item Expand URL, DNS, HTTP, HTTPS, TLS, IP, TLD, and ISP.
  \item Label every part of \texttt{https://blog.example.org:443/notes/net?ch=1}.
  \item Distinguish a URL from a domain name.
  \item Describe the steps of DNS resolution when a name is not cached.
  \item Contrast HTTP with HTTPS on security, default port, and address-bar
    form.
\end{enumerate}

\subsection*{Solutions}
\begingroup\small
\begin{enumerate}[label=\arabic*.]
  \item URL --- Uniform Resource Locator; DNS --- Domain Name System; HTTP
    --- Hypertext Transfer Protocol; HTTPS --- Hypertext Transfer Protocol
    Secure; TLS --- Transport Layer Security; IP --- Internet Protocol; TLD
    --- Top-Level Domain; ISP --- Internet Service Provider.
  \item Scheme \texttt{https}; subdomain \texttt{blog}; domain
    \texttt{example}; TLD \texttt{.org}; port \texttt{443}; path
    \texttt{/notes/net}; query \texttt{?ch=1}.
  \item A URL is the full address of one resource (scheme, host, port,
    path, query). A domain name is only the readable host part inside it.
  \item The browser checks its cache; then the resolver asks the DNS
    hierarchy (root, then TLD, then the authoritative server), which returns
    the IP address; the resolver caches it and hands it to the browser.
  \item HTTP carries pages without encryption over port 80
    (\texttt{http://}); HTTPS is HTTP protected by TLS over port 443
    (\texttt{https://} with a padlock).
\end{enumerate}
\endgroup
\end{document}
